% Follow-on research fragment. Requires the definitions and lemmas of
% moment_rank.tex, including pi_n, c_n, h_n, and Bernoulli 2-adic parity.

\section{The optimal prefix for every pair \texorpdfstring{$(4,q)$}{(4,q)}}
\label{rf:section}

The same orthogonal polynomials settle a second infinite unequal-order
family.  The perturbation now has rank two.  Its determinant involves
three consecutive even Bernoulli numbers, and their elementary
$2$-adic parity proves that this determinant never vanishes.

Let $p=4$, $q=2b$, and $b\geq3$.  In this section write
\[
 M_b^{[4]}=(a_{b+t,k})_{0\leq t,k<b}.
\]
For $b\geq4$, put
\[
 d_b=4b^2-6b-1,\qquad
 u=B_{2b-6},\quad v=B_{2b-4},\quad w=B_{2b-2},
\]
and define
\begin{align}
 \Xi_b={}&4d_b
 +(b-2)(2b-5)\{u+(2b-1)w(4-u)\}\notag\\
 &+(b-1)(2b-3)\{(2b-1)v^2-4(b+1)v\}.
 \label{rf:Xi}
\end{align}

\begin{theorem}[Exact determinant for the entire order-four family]
\label{rf:main}
For every integer $b\geq4$,
\begin{equation}\label{rf:det}
 \boxed{\displaystyle
 \det M_b^{[4]}=-\frac{\Xi_b}
      {6\,2^b(b-1)(2b-4)!}\ne0.}
\end{equation}
Moreover $4\Xi_b$ belongs to $\mathbb Z_{(2)}$ and is odd there;
equivalently, $v_2(\Xi_b)=-2$.  At the remaining unequal-order
case $b=3$,
\begin{equation}\label{rf:smallcase}
 M_3^{[4]}=
 \begin{pmatrix}
 0&3/4&1/4\\5/12&2/3&-1/12\\1&-3/8&1/12
 \end{pmatrix},\qquad
 \det M_3^{[4]}=-\frac{113}{384}.
\end{equation}
Thus the first $q$ centered moments have full formal rank $q$
for every even pair $(p,q)=(4,2b)$ with $q>4$.  Their complete
coefficient determinant is the determinant in
\eqref{rf:det} or \eqref{rf:smallcase}, divided by $(4b-1)!!$.
\end{theorem}

\subsection{A shifted Hankel matrix and its first two inverse rows}

Retain $\mu_n=B_{2n}/(2n)!$ for $n\geq0$, and put $\mu_{-1}=0$.
Define
\[
 H_b^{[-1]}=(\mu_{t+k-1})_{0\leq t,k<b}.
\]
The source formula gives the exact factorization
\begin{equation}\label{rf:factorization}
 M_b^{[4]}=\mathcal R_b(H_b^{[-1]}+U_b)\mathcal E_b,
\end{equation}
where
\[
 (\mathcal R_b)_{tt}=(2b+2t-3)!,\qquad
 (\mathcal E_b)_{kk}=\frac1{(2b-2k)!},
\]
and $U_b$ is zero outside its first two columns, with
\begin{equation}\label{rf:U}
 (U_b)_{tk}=\mu_{t+k-1}
  \frac{(2t+1)!(2b-2k)!}
       {(4-2k)!(2b+2t-3)!},\qquad k=0,1.
\end{equation}
This is precisely where the exclusion of a zero lower binomial
index matters: no third perturbation column is present.

Let $\mathcal L_{-1}(x^j)=\mu_{j-1}$, with
$\mathcal L_{-1}(1)=0$.  The shifted functional from the preceding
section satisfies
\[
 \mathcal L_{-1}(x^2g)=\mathcal N(g).
\]
Thus the explicit monic polynomials $\pi_n$ are orthogonal for the
functional on the right.  The Pad\'e coefficient identities already
proved give
\begin{align}
 \mathcal L_{-1}(x\pi_n)&=c_n,\notag\\
 \mathcal L_{-1}(\pi_n)&=c_n\mathfrak r_n,
 \qquad
 \mathfrak r_n=\frac{(2n-1)(2n)(2n+3)(2n+4)}6.
 \label{rf:two-moments}
\end{align}
The latter formula is valid also at $n=0$, where both sides vanish.
For $n\geq1$ it follows by taking the coefficient of $z^{n-1}$
in $F(z)q_n(z)$: that coefficient equals the coefficient of
$z^{n-1}$ in the explicit numerator $r_n(z)$.  The quotient of
this numerator coefficient and its leading coefficient $c_n$ is
exactly $\mathfrak r_n$.

Set
\[
 n=b-1,\quad m=b-2,\quad
 r_+=\mathfrak r_n,\quad r_-=\mathfrak r_m,
 \quad \delta_b=r_+-r_-
      =\frac{2(4b-3)d_b}{3}.
\]
Both $c_n$ and $c_m$ are nonzero, and $\delta_b>0$ for $b\geq2$.
Replace the last two monomial rows in $H_b^{[-1]}$ by the rows
for $\pi_m,\pi_n$.  This triangular row change has determinant
one, and these rows annihilate columns $2,\ldots,b-1$ by
$\mathcal N$-orthogonality.  Expanding by the first two columns
therefore gives
\begin{equation}\label{rf:Hdet}
 \det H_b^{[-1]}
 =-c_mc_n\delta_b\prod_{j=0}^{b-3}h_j\ne0.
\end{equation}
The two first rows of its inverse are consequently the coefficient
rows of the explicit polynomials
\begin{align}
 V_0(x)&=\frac1{\delta_b}
       \left(\frac{\pi_n(x)}{c_n}-\frac{\pi_m(x)}{c_m}\right),
       \label{rf:inverse0}\\
 V_1(x)&=\frac1{\delta_b}
       \left(-r_-\frac{\pi_n(x)}{c_n}
                       +r_+\frac{\pi_m(x)}{c_m}\right).
       \label{rf:inverse1}
\end{align}
Indeed, the two moments in \eqref{rf:two-moments} give the required
Kronecker values in columns zero and one, and orthogonality kills
all remaining columns.

Combining \eqref{rf:Hdet} with the diagonal factors gives
\begin{equation}\label{rf:base-scaled}
 \det\mathcal R_b\det\mathcal E_b\det H_b^{[-1]}
 =-\frac{d_b}{6\,2^b(b-1)(2b-4)!}\qquad(b\geq2).
\end{equation}
This factorial simplification can be checked inductively: after
division by $d_b$, the value at $b=2$ is $-1/24$, and the quotient
of successive values is $1/[4(b-1)(2b-5)]$.

\subsection{The four perturbation contractions}

For $j=n,m$ and $k=0,1$, put
\[
 s_{j,k}=\frac1{c_j}\sum_{t=0}^j[x^t]\pi_j(x)\,(U_b)_{tk},
 \qquad s_{+,k}=s_{n,k},\quad s_{-,k}=s_{m,k}.
\]
The explicit factorial coefficients of $\pi_j$ simplify these
contractions to the following finite Bernoulli sums:
\begin{align}
 s_{+,0}&=\frac{2b-2}{12}
 \sum_{\substack{0\leq k\leq2b-4\\k\ {\rm even}}}
 \binom{2b-3}{k}(k+3)(k+2b)(k+2b+1)B_k,\notag\\
 s_{+,1}&=\frac1{2b(2b-1)^2}
 \sum_{\substack{0\leq k\leq2b-2\\k\ {\rm even}}}
 \binom{2b-1}{k}(k+1)(k+2b-2)(k+2b-1)B_k,\notag\\
 s_{-,0}&=\frac{(2b)(2b-1)(2b-4)}{12}
 \sum_{\substack{0\leq k\leq2b-6\\k\ {\rm even}}}
 \binom{2b-5}{k}(k+3)B_k,\notag\\
 s_{-,1}&=\frac1{2b-3}
 \sum_{\substack{0\leq k\leq2b-4\\k\ {\rm even}}}
 \binom{2b-3}{k}(k+1)B_k.
 \label{rf:finite-sums}
\end{align}

For completeness, the polynomially weighted Bernoulli identities
needed to evaluate these sums are as follows.  If $N\geq3$ is odd,
\begin{equation}\label{rf:linear-sum}
 \sum_{k\ {\rm even}}\binom Nk(k+\alpha)B_k
 =N\left(\frac{\alpha+1}{2}-B_{N-1}\right).
\end{equation}
If $N\geq5$ is odd and $P$ is a monic cubic, then
\begin{align}
 \sum_{k\ {\rm even}}\binom NkP(k)B_k
 ={}&\frac N2P(1)
   +N\{P(N-1)-P(N)\}B_{N-1}\notag\\
 &-N(N-1)(N-2)B_{N-3}.
 \label{rf:cubic-sum}
\end{align}
Here and in \eqref{rf:linear-sum} the sum runs over $0\leq k\leq N$.
To prove the identities, first include the $k=1$ term, then expand
$P(N-y)$ in the falling factorials of $y=N-k$.  For $0\leq j\leq3$,
\[
 \sum_{k=0}^N(N-k)_{\underline j}\binom NkB_k
 =N_{\underline j}B_{N-j}(1).
\]
When $N\geq5$ is odd, only $j=1,3$ survive.  The falling-factorial
coefficient at $j=1$ is $P(N-1)-P(N)$, and the coefficient at
$j=3$ is $-1$.  Removing the $k=1$ term adds $NP(1)/2$.
The linear version is obtained in exactly the same way.  Thus these
are polynomial identities, not conjectural Bernoulli reductions.

Applying them to \eqref{rf:finite-sums}, for $b\geq4$ one obtains
\begin{align}
 s_{-,0}&=r_-(1-u/2),& s_{-,1}&=1-v,\notag\\
 s_{+,0}&=r_+
 -\frac{(2b-3)(2b-2)(2b-1)(4b-3)}3v\notag\\
 &\hspace{9mm}
 -\frac{(2b-3)(2b-2)(2b-4)(2b-5)}{12}u,\notag\\
 s_{+,1}&=1-\frac{2(4b-3)}b w
             -\frac{(b-1)(2b-3)}{b(2b-1)}v.
 \label{rf:four-contractions}
\end{align}
The separate treatment of $b=3$ avoids applying
\eqref{rf:linear-sum} at $N=1$, where that formula is not valid.

The rank-two determinant lemma and
\eqref{rf:inverse0}--\eqref{rf:inverse1} now give
\begin{equation}\label{rf:small-det}
 \frac{\det(H_b^{[-1]}+U_b)}{\det H_b^{[-1]}}
 =\frac{(r_++s_{+,0})(1+s_{-,1})
             -(r_-+s_{-,0})(1+s_{+,1})}{\delta_b}
 =\frac{\Xi_b}{d_b}.
\end{equation}
The last equality is direct polynomial simplification of the four
displayed contractions.  Together with \eqref{rf:base-scaled},
it proves the determinant formula \eqref{rf:det}.

\subsection{A uniform arithmetic proof of nonvanishing}

The Bernoulli parity argument in the order-two section proves that
\[
 \widehat u=2u,\qquad\widehat v=2v,\qquad\widehat w=2w
\]
are odd elements of $\mathbb Z_{(2)}$, since $b\geq4$ makes all
three Bernoulli indices positive and even.  Multiplication of
\eqref{rf:Xi} by four yields
\begin{align}
 4\Xi_b={}&16d_b
 +(b-2)(2b-5)
 \{2\widehat u+(2b-1)\widehat w(8-\widehat u)\}\notag\\
 &+(b-1)(2b-3)
 \{(2b-1)\widehat v^2-8(b+1)\widehat v\}.
 \label{rf:parity}
\end{align}
Modulo two, the first bracket contributes $b$ and the second
contributes $b-1$.  Hence
\[
 4\Xi_b\equiv b+(b-1)\equiv1\pmod2.
\]
Therefore $v_2(\Xi_b)=-2$ and $\Xi_b\ne0$ for every $b\geq4$.
This proves nonvanishing uniformly, without a real-sign assertion
or any assumption about Bernoulli numerators beyond the elementary
parity already established.

\paragraph{Updated scope of the initial-prefix problem.}
The optimal first-$q$ prefix is now proved for all even pairs with
$p=2$, all even pairs with $p=4$, and the diagonal $p=q$ inherited
from the source.  The next section settles $p=6$; the final remaining unequal-order
range in this article is $8\leq p<q$.  The rank-two calculation suggests investigating the
$2$-adic structure of the analogous higher-rank Bernoulli
contractions, but no general nonvanishing theorem for those
contractions is assumed here.

% Optional compact corollary of the explicit (4,6) source matrix.
\subsection{A convergent harmonic--polygamma identity at orders four and six}

The exceptional matrix \eqref{rf:smallcase} also gives a concrete
ordinary convergent identity.  Let $x_n=n+\tfrac12$ and
$A_s(n)=\zeta(s)-H_n^{(s)}$, and put
\[
 \begin{split}
 W_{4,6}(x)={}&67+15516x^2+324320x^4+1434496x^6\\
              &+1601280x^8+371712x^{10}.
 \end{split}
\]
Then
\begin{equation}\label{rf:ordinary-46}
 \boxed{\displaystyle
 \sum_{n=0}^{\infty}\left\{
 W_{4,6}(x_n)A_4(n)A_6(n)
 -\frac{123904x_n^2+316928}{5}\right\}=-21696.}
\end{equation}
The summand is $O(n^{-2})$, so the convergence is absolute.
One may equivalently write the tail product as
\[
 A_4(n)A_6(n)
 =\frac{\psi^{(3)}(n+1)\psi^{(5)}(n+1)}{720}.
\]

\begin{proof}
Direct multiplication of the exact matrix in
\eqref{rf:smallcase} gives
\[
 (-7,45,66)M_3^{[4]}=(339/4,0,0),
 \qquad -7C_3+45C_4+66C_5=-\frac{339}{8}.
\]
Set $S(y)=512(-7y^7+45y^9+66y^{11})$.  The inverse Bernoulli
transform identifies the polynomial of the resulting centered
moment combination as
\[
 S(x+\tfrac12)-S(x-\tfrac12)=W_{4,6}(x).
\]
The centered Hurwitz expansions give
\[
 A_4(n)A_6(n)=\frac1{15}x_n^{-8}
              -\frac7{60}x_n^{-10}+O(x_n^{-12}).
\]
Hence the polynomial in nonnegative powers of $x_n$ in its product
with $W_{4,6}$ is exactly
$(123904x_n^2+316928)/5$.  Subtracting it leaves an absolutely
summable remainder.  The removed polynomial has continued sum zero;
equivalently, its partial sum has zero constant term in the integer cutoff, by the
centered even-power argument used in
Corollary~\ref{cor:ordinary-two}.  The ordinary sum therefore equals
the finite-part moment combination, namely
$512(-339/8)=-21696$.
\end{proof}
